%% ma: L11-B17-0018
%% lop: 11
%% chuong: 5
%% bai: 17
%% dang: 11.17.3
%% loai: TL
%% muc: VD
%% dap_an: "-"
%% nguon: "tai-lieu-hoc-tap-mon-toan-lop-11-theo-chuong-trinh-giao-duc-pho-thong-2018.pdf (Lớp Toán Thầy Duy, Phạm Lê Duy, THPT Chu Văn An), Chương 3 Bài 3, Dạng 3, Ví dụ 18 (tr.393-398)"
%% kien_thuc: "Lớp 11 Bài 17 (hàm số liên tục)"
%% trang_thai: cho-duyet
%% ly_do: "Dạng mới đề xuất (11.17.x), chờ giáo viên duyệt dạng; tài liệu gốc không có lời giải mẫu, đáp án do kho tự giải và kiểm bằng sympy."
\begin{ex}
Chứng minh rằng phương trình $\dfrac1{\cos x}-\dfrac1{\sin x}=m$ luôn có nghiệm với mọi giá trị của tham số $m$.
\loigiai{\sloppy Hàm $f(x)=\dfrac1{\cos x}-\dfrac1{\sin x}$ liên tục trên $\left(0;\dfrac\pi2\right)$. Khi $x\to0^+$: $\dfrac1{\cos x}\to1$, $\dfrac1{\sin x}\to+\infty$ nên $f(x)\to-\infty$; khi $x\to\dfrac{\pi}2^-$: $\dfrac1{\cos x}\to+\infty$, $\dfrac1{\sin x}\to1$ nên $f(x)\to+\infty$. Với mỗi $m$ tồn tại $x_1,x_2\in\left(0;\dfrac\pi2\right)$ sao cho $f(x_1)<m<f(x_2)$. Theo định lí giá trị trung gian, $f(x)=m$ có nghiệm nằm giữa $x_1$ và $x_2$.}
\end{ex}
